% Compléments Bloc 2 — zooms techniques

\begin{frame}{C2.1 — S-123 : comparer les deux architectures en un regard}
\context{B2}{C2.1 · S-123 · Zoom preuve}
\scriptsize
\begin{tabularx}{\linewidth}{>{\bfseries}p{3.0cm}X X}
\toprule
Critère & Modèle direct S-123 & Méta-modèle retenu \\
\midrule
Structure & une table par concept métier & concepts génériques du Feature Catalogue \\
Volumétrie du schéma & 73 tables pour une version partielle & 18 tables principales \\
Évolution & migrations fréquentes à chaque version & chargement de nouvelles définitions \\
Réutilisation & fortement lié à S-123 & compatible avec plusieurs S-1xx \\
Maintenance & risque d'obsolescence rapide & abstraction plus stable \\
\bottomrule
\end{tabularx}
\begin{block}{Leçon d'architecture}
Une bonne architecture ne copie pas mécaniquement une norme : elle choisit le niveau d'abstraction adapté à son évolution.
\end{block}
\end{frame}

\begin{frame}{C2.2 — piv2glz : rendre l'intégrité du code visible}
\context{B2}{C2.2 · piv2glz · Zoom preuve}
\begin{columns}[T]
\begin{column}{0.47\linewidth}
\begin{block}{Tests}
\begin{tightitems}
  \item Fonctions isolées et cas nominaux.
  \item Erreurs de système de fichiers provoquées avec \texttt{monkeypatch}.
  \item Scénarios CLI réels avec \texttt{subprocess}.
  \item Fichiers temporaires via \texttt{tmp\_path}.
\end{tightitems}
\end{block}
\end{column}
\begin{column}{0.49\linewidth}
\begin{block}{CI}
\begin{tightitems}
  \item Environnement Python propre.
  \item Pytest exécuté avant la construction.
  \item PyInstaller seulement après succès.
  \item Artefacts Windows récupérables depuis GitLab.
\end{tightitems}
\end{block}
\end{column}
\end{columns}
\begin{block}{Pourquoi cette preuve est forte}
Le périmètre est petit, mais le mécanisme qualité est entièrement reproductible et automatisé.
\end{block}
\end{frame}

\begin{frame}{C2.3 — Brise-glace : le front-end est piloté par le scénario utilisateur}
\context{B2}{C2.3 · Brise-glace · Zoom preuve}
\begin{columns}[T]
\begin{column}{0.46\linewidth}
\begin{block}{Contraintes d'interface}
\begin{tightitems}
  \item États différents selon la phase de jeu.
  \item Rôles joueur / maître du jeu.
  \item Interaction par glisser-déposer.
  \item Utilisation ordinateur et mobile.
\end{tightitems}
\end{block}
\end{column}
\begin{column}{0.50\linewidth}
\begin{block}{Réponse technique}
\begin{tightitems}
  \item React + TypeScript.
  \item État partagé synchronisé en temps réel.
  \item Supabase pour persistance et événements.
  \item Validation fonctionnelle par scénarios de partie.
\end{tightitems}
\end{block}
\end{column}
\end{columns}
\centering
\repofigure[0.27\textheight]{figures/brise_glace/sequence_partie.png}
\end{frame}

\begin{frame}{C2.5 — Distinguer ce qui existe de l'architecture de passage à l'échelle}
\context{B2}{C2.5 · Data/IA · Zoom preuve}
\begin{columns}[T]
\begin{column}{0.47\linewidth}
\begin{block}{Réalisé}
\begin{tightitems}
  \item Préparation et exploration des données.
  \item Comparaison de modèles.
  \item Explicabilité.
  \item FastAPI et historique des prédictions.
\end{tightitems}
\end{block}
\end{column}
\begin{column}{0.49\linewidth}
\begin{block}{Si les contraintes changent}
\begin{tightitems}
  \item Ingestion séparée.
  \item Stockage brut et exploitation distincts.
  \item Traitement distribué si la volumétrie le justifie.
  \item Cycle de vie et surveillance du modèle.
\end{tightitems}
\end{block}
\end{column}
\end{columns}
\begin{block}{Positionnement}
Je ne transforme pas rétrospectivement le projet en Big Data : j'explicite comment l'architecture évoluerait si le besoin l'exigeait.
\end{block}
\end{frame}
